\documentclass[11pt]{article}
\usepackage[margin=0.82in]{geometry}
\usepackage[T1]{fontenc}
\usepackage{lmodern}
\usepackage{microtype}
\usepackage{amsmath,amssymb,amsthm,mathtools}
\usepackage{booktabs,array,enumitem}
\usepackage{xcolor}
\usepackage{fancyhdr}
\usepackage{hyperref}

\setlength{\parindent}{0pt}
\setlength{\parskip}{0.55em}
\setlist[itemize]{leftmargin=1.4em,itemsep=0.25em,topsep=0.25em}
\setlist[enumerate]{leftmargin=1.6em,itemsep=0.25em,topsep=0.25em}

\newtheorem{theorem}{Theorem}[section]
\newtheorem{proposition}[theorem]{Proposition}
\newtheorem{corollary}[theorem]{Corollary}
\theoremstyle{definition}
\newtheorem{definition}[theorem]{Definition}

\pagestyle{fancy}
\fancyhf{}
\fancyhead[L]{\footnotesize\scshape Language Mechanics}
\fancyhead[R]{\footnotesize\scshape Minimum Root Algebra v0.2}
\fancyfoot[C]{\footnotesize\thepage}
\renewcommand{\headrulewidth}{0.35pt}
\setlength{\headheight}{14pt}
\fancypagestyle{first}{\fancyhf{}\fancyfoot[C]{\footnotesize\thepage}\renewcommand{\headrulewidth}{0pt}}

\newcommand{\R}{\mathbb R}
\newcommand{\Mtwo}{M_2(\R)}
\newcommand{\Alg}{\operatorname{Alg}_{\R}}
\newcommand{\Span}{\operatorname{span}_{\R}}
\newcommand{\diag}{\operatorname{diag}}
\newcommand{\Dih}{\operatorname{Dih}}

\begin{document}
\thispagestyle{first}

\begin{center}
{\Large\bfseries 010 Minimum Root Algebra}\par
\vspace{0.08in}
{\large A Two-Generator Kernel for $M_2(\R)$}\par
\vspace{0.16in}
{Anthony Vito Coppa}\par
\vspace{0.05in}
{\small Language Mechanics -- Formal Form 010}\par
{\small Standalone edition v0.2 -- 29 August 2026}\par
\end{center}
\vspace{0.08in}
\hrule
\vspace{0.12in}

\textbf{Status.} Carrier choices and definitions are \textsc{Declared}. The generator-minimum theorem, unique-coordinate corollary, multiplication identities, Fibonacci recurrence, golden null/eigenframe result, and reversal identities are \textsc{Derived} from the displayed matrices. Independent re-audit of this standalone edition remains \textsc{Open}.

\textbf{Scope.} This paper begins with two declared real $2\times2$ matrices. It proves the minimum generator cardinality for the full unital real algebra $\Mtwo$ and derives the exact structures carried by that pair. It does not claim that the displayed matrix $Q$ is uniquely selected from every possible transporter class. Here \emph{minimum} means only generator cardinality inside $\Mtwo$.

\textbf{Question.} What is the minimum number of real $2\times2$ generators required to produce the full noncommutative algebra $\Mtwo$, and what exact normal form follows from one explicit minimum pair?

\section{The carrier and its two generators}

Declare
\[
Q=\begin{pmatrix}1&1\\[2pt]1&0\end{pmatrix},
\qquad
C_1=\begin{pmatrix}0&-1\\[2pt]1&0\end{pmatrix},
\qquad
\Omega:=-C_1=\begin{pmatrix}0&1\\[2pt]-1&0\end{pmatrix}.
\]
Then direct multiplication gives
\[
C_1^2=-I,\qquad C_1^4=I,\qquad Q^2=Q+I,\qquad Q^{-1}=Q-I.
\]
Define the order defect
\[
G:=[C_1,Q]=C_1Q-QC_1=Q\Omega-\Omega Q
 =\begin{pmatrix}-2&1\\[2pt]1&2\end{pmatrix}.
\]
Again directly,
\[
G^2=5I,\qquad G^{\mathsf T}=G.
\]

\section{Generator minimum and complete local algebra}

\begin{theorem}[Two generators are cardinality-minimal]
The ordered set
\[
\mathcal B=(I,Q,\Omega,G)
\]
is a basis of $\Mtwo$. Consequently
\[
\Alg(Q,C_1)=\Mtwo.
\]
No single real $2\times2$ matrix generates $\Mtwo$ as a unital real algebra.
\end{theorem}

\begin{proof}
Vectorize the four matrices in the coordinate order $(a,b,c,d)$ for $\left(\begin{smallmatrix}a&b\\c&d\end{smallmatrix}\right)$. The corresponding basis matrix is
\[
B=
\begin{pmatrix}
1&1&0&-2\\
0&1&1&1\\
0&1&-1&1\\
1&0&0&2
\end{pmatrix},
\qquad \det B=-10\neq0.
\]
Hence $I,Q,\Omega,G$ are linearly independent and therefore span the four-dimensional vector space $\Mtwo$. Each lies in the unital algebra generated by $Q$ and $C_1$, because $\Omega=-C_1$ and $G=[C_1,Q]$.

For any single $A\in\Mtwo$, Cayley--Hamilton reduces every polynomial in $A$ to $xI+yA$. Thus $\R[A]$ is commutative and has dimension at most two. It cannot equal the four-dimensional noncommutative algebra $\Mtwo$. Therefore two generators are necessary and sufficient.\qedhere
\end{proof}

\begin{center}
\fbox{\begin{minipage}{0.88\linewidth}\small
\textbf{Minimum statement.} Two real matrix generators are necessary and sufficient for this carrier. ``Minimum'' here means generator cardinality in the unital real algebra $M_2(\R)$; it does not assert a universal minimum over unrelated carrier classes.
\end{minipage}}
\end{center}

\section{Unique four-coordinate normal form}

\begin{corollary}
Every
\[
A=\begin{pmatrix}a&b\\c&d\end{pmatrix}\in\Mtwo
\]
has exactly one representation
\[
A=x_0I+x_1Q+x_2\Omega+x_3G,
\]
where
\[
 x_0=\frac{2a+3d-b-c}{5},\qquad
 x_1=\frac{a-d+2b+2c}{5},
\]
\[
 x_2=\frac{b-c}{2},\qquad
 x_3=\frac{-2a+2d+b+c}{10}.
\]
\end{corollary}

\begin{proof}
Solving the linear system $B(x_0,x_1,x_2,x_3)^{\mathsf T}=(a,b,c,d)^{\mathsf T}$ gives the displayed coordinates. Uniqueness follows because $\det B\neq0$.
\end{proof}

This theorem supplies a matrix-value normal form only. Word provenance, evaluated matrix value, and occurrence receipt remain distinct objects.

\section{Exact multiplication table}

Multiplication in the basis $\mathcal B$ is
\[
\begin{array}{c|cccc}
\cdot & I & Q & \Omega & G\\ \hline
I      & I & Q & \Omega & G\\
Q      & Q & I+Q & (\Omega+G)/2 & (G+5\Omega)/2\\
\Omega & \Omega & (\Omega-G)/2 & -I & 2Q-I\\
G      & G & (G-5\Omega)/2 & I-2Q & 5I
\end{array}
\]
(row times column). Each entry follows by direct multiplication of the displayed matrices. By bilinearity, the table reduces every finite word in $\{Q,C_1,Q^{-1},C_1^{-1}\}$ to the unique normal form of Section 3.

\textbf{Example.} Since $C_1=-\Omega$ and $C_1^{-1}=\Omega$,
\[
C_1QC_1^{-1}=-\Omega Q\Omega.
\]
Using the table,
\[
\Omega Q\Omega
=\frac{\Omega-G}{2}\,\Omega
=\frac{-I-(I-2Q)}{2}
=Q-I=Q^{-1},
\]
so
\[
\boxed{C_1QC_1^{-1}=-Q^{-1}.}
\]

\section{Fibonacci carriage}

The quadratic identity $Q^2=Q+I$ gives the complete recurrence.

\begin{theorem}
For $n\ge1$,
\[
\boxed{Q^n=F_nQ+F_{n-1}I},
\]
where $F_0=0$, $F_1=1$, and $F_{n+1}=F_n+F_{n-1}$.
\end{theorem}

\begin{proof}
The case $n=1$ is immediate. If the formula holds at $n$, then
\[
Q^{n+1}=F_nQ^2+F_{n-1}Q
=F_n(Q+I)+F_{n-1}Q
=F_{n+1}Q+F_nI.
\]
\end{proof}

The recurrence follows from the declared carrier. It is not an independent selector for $Q$.

\section{Golden eigenframe and null structure}

Let
\[
\Phi=\frac{1+\sqrt5}{2},\qquad
\Psi=\frac{1-\sqrt5}{2},
\]
so $\Phi+\Psi=1$, $\Phi\Psi=-1$, and $\Phi-\Psi=\sqrt5$. Define
\[
v_\Phi=\binom{\Phi}{1},\qquad
v_\Psi=\binom{\Psi}{1}.
\]

\begin{proposition}[Golden eigenlines are exactly the null lines]
One has
\[
Qv_\Phi=\Phi v_\Phi,\qquad Qv_\Psi=\Psi v_\Psi.
\]
For
\[
\Delta_G(x,y):=(x,y)G\binom{x}{y}
=-2x^2+2xy+2y^2,
\]
its nonzero projective null locus is exactly
\[
\Delta_G(v)=0
\iff
[v]\in\{[v_\Phi],[v_\Psi]\}.
\]
Moreover $G$ has signature $(1,1)$ and
\[
v_\Phi^{\mathsf T}\Omega v_\Psi=\sqrt5.
\]
With
\[
H:=\frac{G}{\sqrt5},
\]
one has
\[
H^2=I,\qquad H^{\mathsf T}H=I,\qquad \det H=-1.
\]
Thus $H$ is a Euclidean reflection constructed from the generator order defect.
\end{proposition}

\begin{proof}
The eigenvector equations reduce to $r^2=r+1$, whose roots are $\Phi$ and $\Psi$. Also
\[
\Delta_G(x,y)=-2(x^2-xy-y^2).
\]
For a nonzero null vector, $y\neq0$; setting $r=x/y$ gives $r^2-r-1=0$, hence $r\in\{\Phi,\Psi\}$. Since $\det G=-5<0$, the symmetric form has signature $(1,1)$. Finally,
\[
v_\Phi^{\mathsf T}\Omega v_\Psi=\Phi-\Psi=\sqrt5.
\]
Because $G^{\mathsf T}=G$ and $G^2=5I$, we have $H^{\mathsf T}H=H^2=I$, while $\det H=\det(G)/5=-1$.
\end{proof}

\section{Exact reversal identities}

\begin{proposition}[Finite and infinite reversal]
The reflection $H=G/\sqrt5$ satisfies
\[
\boxed{HC_1H=C_1^{-1}},\qquad
\boxed{HQH=-Q^{-1}}.
\]
Consequently, with $B:=Q^2$,
\[
HBH=B^{-1}.
\]
Since $B$ has infinite order and $H\notin\langle B\rangle$, the generated group obeys the infinite-dihedral presentation
\[
\langle B,H\mid H^2=I,\ HBH=B^{-1}\rangle\cong D_\infty.
\]
\end{proposition}

\begin{proof}
From the multiplication table,
\[
G\Omega=I-2Q,
\qquad
QG=\frac{G+5\Omega}{2}.
\]
Hence
\[
G\Omega G=(I-2Q)G
=G-(G+5\Omega)
=-5\Omega.
\]
Because $C_1=-\Omega$,
\[
HC_1H
=-\frac{G\Omega G}{5}
=\Omega=C_1^{-1}.
\]
Likewise,
\[
GQ=\frac{G-5\Omega}{2},\qquad \Omega G=2Q-I,
\]
so
\[
GQG
=\frac{G^2-5\Omega G}{2}
=\frac{5I-5(2Q-I)}{2}
=5(I-Q)
=-5Q^{-1}.
\]
Therefore $HQH=-Q^{-1}$. Squaring the $Q$ relation gives $HBH=B^{-1}$.

The eigenvalues of $B=Q^2$ are $\Phi^2$ and $\Phi^{-2}$, so $B$ has infinite order. Also $\det B=1$ while $\det H=-1$, hence $H\notin\langle B\rangle$. The displayed relations therefore give the standard infinite-dihedral group.\qedhere
\end{proof}

\section{Exact boundary}

This form proves only the following:
\begin{itemize}
\item two generators are minimal by cardinality for the unital real algebra $M_2(\R)$;
\item $Q$ and $C_1$ generate all of $M_2(\R)$;
\item every matrix has one four-coordinate normal form in $(I,Q,\Omega,G)$;
\item the displayed multiplication table reduces every finite generator word;
\item Fibonacci powers follow from $Q^2=Q+I$;
\item the golden eigenlines of $Q$ are exactly the null lines of $G$;
\item $H=G/\sqrt5$ is a reflection satisfying the displayed finite and infinite reversal identities.
\end{itemize}

It does not prove that $Q$ is the unique or universal transporter outside a separately declared selection class. It supplies no automatic identity with physical fields, spacetime geometry, gauge structure, spin, atomic structure, biological memory, semantic meaning, an independent eight-address carrier, or a twenty-eight-seam construction. Any such use requires its own typed construction.

\begin{center}
\fbox{\begin{minipage}{0.88\linewidth}\small
\textbf{Final form.}
\[
(Q,C_1)\longrightarrow G=[C_1,Q]\longrightarrow\{I,Q,\Omega,G\}=M_2(\R)
\longrightarrow\text{unique normal form}.
\]
The carrier is complete. Everything further is construction.
\end{minipage}}
\end{center}

\section*{Source record}
The minimum-root algebra family was first consolidated in August 2026; the source edition titled \emph{Minimum Root Algebra v2} was revised on 18 August 2026. This standalone Language Mechanics edition uses that source only as provenance. Every mathematical claim used here is stated and proved or directly computed in the present form.

\end{document}
